% !TEX root = ../algorithms.tex

\section{Линейное программирование и метод внутренней точки}

Рассматривается стандартная пара задач линейного программирования
\[
\cP:\quad
\max \; c^\top x
\quad \text{при условиях} \quad
Ax \le b,
\]
\[
\cD:\quad
\min \; b^\top y
\quad \text{при условиях} \quad
A^\top y = c,\; y \ge 0,
\]
где \(A \in \Q^{m \times n}\), \(b \in \Q^m\), \(c \in \Q^n\).

\begin{Proposition}{Слабая двойственность}{weak-duality}
	Для любых допустимых решений \(x\) прямой задачи и \(y\) двойственной задачи выполнено
	\[
	c^\top x \le b^\top y.
	\]
\end{Proposition}

\begin{proof}
	Так как \(A^\top y = c\), имеем
	\[
	c^\top x = (A^\top y)^\top x = y^\top A x.
	\]
	Далее, из \(Ax \le b\) и \(y \ge 0\) следует
	\[
	y^\top A x \le y^\top b = b^\top y.
	\]
	Следовательно,
	\[
	c^\top x \le b^\top y.
	\]
\end{proof}

\begin{Remark}{Сведение ЛП к задаче выполнимости}{}
	Предположим, что имеется алгоритм, который за полиномиальное время решает задачу выполнимости системы линейных неравенств:
	либо находит решение, либо сообщает, что решений нет.
	
	Тогда задачу линейного программирования можно свести к выполнимости одной системы линейных ограничений. Действительно,
	\(x\) и \(y\) являются оптимальными решениями задач \(\cP\) и \(\cD\) тогда и только тогда, когда выполнена система
	\[
	Ax \le b,\qquad A^\top y = c,\qquad y \ge 0,\qquad b^\top y - c^\top x \le 0.
	\]
	По слабой двойственности для любых допустимых \(x\) и \(y\) всегда выполнено
	\[
	b^\top y - c^\top x \ge 0.
	\]
	Следовательно, добавленное неравенство \(b^\top y - c^\top x \le 0\) эквивалентно равенству
	\[
	b^\top y = c^\top x,
	\]
	то есть условию оптимальности. Таким образом, нахождение оптимального решения сводится к нахождению решения системы линейных ограничений.
	
	Равенство \(A^\top y = c\) при желании можно заменить на пару систем неравенств
	\[
	A^\top y \le c,
	\qquad
	-A^\top y \le -c,
	\]
	так что всё сводится именно к проверке выполнимости системы линейных неравенств.
\end{Remark}

\begin{Proposition}{Задачи \(\cP\) и \(\cD\) образуют двойственную пару}{primal-dual-pair}
	Задача \(\cD\) является двойственной к задаче \(\cP\), а задача \(\cP\) является двойственной к задаче \(\cD\).
\end{Proposition}

\begin{proof}
	Пусть \(x\) — допустимое решение задачи \(\cP\), а \(y\) удовлетворяет условиям
	\[
	A^\top y = c,\qquad y \ge 0.
	\]
	Тогда
	\[
	c^\top x = (A^\top y)^\top x = y^\top A x.
	\]
	Так как \(Ax \le b\) и \(y \ge 0\), имеем
	\[
	y^\top A x \le y^\top b = b^\top y.
	\]
	Следовательно,
	\[
	c^\top x \le b^\top y.
	\]
	
	Итак, любой вектор \(y\), удовлетворяющий условиям
	\[
	A^\top y = c,\qquad y \ge 0,
	\]
	задаёт верхнюю оценку \(b^\top y\) на значение целевой функции прямой задачи.
	Поэтому среди всех таких оценок естественно выбирать наименьшую, то есть решать задачу
	\[
	\min \; b^\top y
	\quad \text{при условиях} \quad
	A^\top y = c,\; y \ge 0.
	\]
	Это и есть задача \(\cD\).
	
	Аналогично, начиная с задачи \(\cD\), получаем, что двойственная к ней задача имеет вид
	\[
	\max \; c^\top x
	\quad \text{при условиях} \quad
	Ax \le b,
	\]
	то есть совпадает с \(\cP\).
\end{proof}

\subsection{Длина входа и полиномиальность}

\begin{Definition}{Длина входа}{input-length}
	Пусть все элементы матрицы \(A\), вектора \(b\) и вектора \(c\) являются рациональными числами,
	каждое из которых кодируется не более чем \(w\) битами.
	Тогда длиной входа задачи линейного программирования называется величина
	\[
	L = (mn + m + n)\, w = \Theta(m n w).
	\]
\end{Definition}

\begin{Definition}{Полиномиальный алгоритм}{polytime}
	Алгоритм решения задачи линейного программирования называется \emph{полиномиальным},
	если его время работы оценивается величиной
	\[
	O\!\bigl(L^{\,O(1)}\bigr).
	\]
\end{Definition}

\subsection{Метод внутренней точки}

\begin{Definition}{Внутренняя точка}{interior-point}
	Точка \(x_0\) называется \emph{внутренней допустимой точкой} для системы \(Ax \le b\),
	если
	\[
	A x_0 < b
	\]
	покомпонентно.
\end{Definition}

\begin{Remark}{Как найти внутреннюю точку}{find-interior}
	Для поиска внутренней точки удобно рассмотреть вспомогательную задачу
	\[
	\max \; t
	\quad \text{при условии} \quad
	A x + t \1 \le b,
	\]
	где \(\1 = (1,\dots,1)^\top \in \R^m\).
	
	Если оптимальное значение \(t^\ast\) положительно, то соответствующая точка \(x^\ast\) удовлетворяет
	\[
	A x^\ast < b,
	\]
	то есть является внутренней точкой исходной системы.
	
	У вспомогательной задачи всегда есть очевидная внутренняя точка, например
	\[
	x = 0,\qquad t = - \max_{1 \le i \le m} |b_i| - 1,
	\]
	поскольку тогда \(t < b_i\) для всех \(i\), а значит
	\[
	A x + t \1 = t \1 < b.
	\]
\end{Remark}

\subsubsection*{Барьерная задача}

Пусть строки матрицы \(A\) имеют вид \(a_1^\top,\dots,a_m^\top\), а
\[
s_i(x) := b_i - a_i^\top x
\]
обозначают зазоры до ограничений.

\begin{Definition}{Логарифмический барьер}{log-barrier}
	Для \(x\), лежащих во внутренности многогранника \(Ax < b\), положим
	\[
	g(x) := \sum_{i=1}^m \ln s_i(x)
	= \sum_{i=1}^m \ln \bigl(b_i - a_i^\top x\bigr).
	\]
	Для параметра \(w > 0\) рассмотрим барьерную задачу
	\[
	\max \; \bigl(w\, c^\top x + g(x)\bigr)
	\quad \text{при условии} \quad
	A x \leq b.
	\]
	Обозначим через \(x_w^\ast\) одно из её решений.
\end{Definition}

\begin{Proposition}{Центральный путь}{}
	Семейство точек \(\{x_w^\ast\}_{w>0}\) образует \emph{центральный путь}. Кроме того,
	\[
	\argmax_{Ax<b} \bigl(w\, c^\top x + g(x)\bigr)
	=
	\argmax_{Ax<b} \left(c^\top x + \frac{1}{w} g(x)\right).
	\]
	Поэтому при \(w \to +\infty\) влияние барьерного слагаемого ослабевает, и точки \(x_w^\ast\)
	приближаются к оптимальному решению исходной задачи линейного программирования.
\end{Proposition}

\begin{Remark}{Геометрический смысл}{geometry}
	Барьерная функция \(g(x)\) стремится к \(-\infty\) при приближении к границе многогранника \(Ax \le b\).
	Поэтому максимизатор \(x_w^\ast\) остаётся во внутренности допустимой области.
	При увеличении \(w\) эти точки движутся внутри многогранника в сторону оптимального решения исходной задачи.
\end{Remark}

\subsubsection*{Идея алгоритма}

Пусть
\[
f_k(x) := w_k\, c^\top x + g(x),
\]
где \(0 < w_0 < w_1 < w_2 < \cdots\).
Идея метода состоит в том, чтобы выбирать параметры \(w_k\) возрастающими
и поддерживать точку \(x_k\), которая достаточно близка к центральной точке
\[
x_k^\ast := x_{w_k}^\ast.
\]
После этого из \(x_k\) выполняется один шаг Ньютона, который перемещает точку ближе к \(x_{k+1}^\ast\).

\subsubsection*{Шаг Ньютона}

\begin{Remark}{Обозначения \(\nabla f\) и \(\nabla^2 f\)}{nabla-notation}
	Пусть \(f(x_1,\dots,x_n)\) — функция многих переменных.
	
	Символ
	\[
	\nabla f(x)
	\]
	обозначает \emph{градиент} функции \(f\) в точке \(x\), то есть столбец её первых частных производных:
	\[
	\nabla f(x)=
	\begin{pmatrix}
		\frac{\partial f}{\partial x_1}(x)\\
		\vdots\\
		\frac{\partial f}{\partial x_n}(x)
	\end{pmatrix}.
	\]
	
	Символ
	\[
	\nabla^2 f(x)
	\]
	обозначает \emph{гессиан} функции \(f\) в точке \(x\), то есть матрицу её вторых частных производных:
	\[
	\nabla^2 f(x)=
	\left(
	\frac{\partial^2 f}{\partial x_i \partial x_j}(x)
	\right)_{i,j=1}^n.
	\]
\end{Remark}

\begin{Proposition}{Градиент и гессиан барьерной функции}{barrier-derivatives}
	Для функции
	\[
	g(x) = \sum_{i=1}^m \ln \bigl(b_i - a_i^\top x\bigr)
	\]
	справедливы формулы
	\[
	\nabla g(x) = - \sum_{i=1}^m \frac{a_i}{s_i(x)},
	\qquad
	\nabla^2 g(x) = - \sum_{i=1}^m \frac{a_i a_i^\top}{s_i(x)^2}.
	\]
	Следовательно,
	\[
	\nabla f_k(x) = w_k c - \sum_{i=1}^m \frac{a_i}{s_i(x)},
	\qquad
	\nabla^2 f_k(x) = - \sum_{i=1}^m \frac{a_i a_i^\top}{s_i(x)^2}.
	\]
\end{Proposition}

\begin{proof}
	Для каждого \(i\) имеем
	\[
	\nabla \ln\bigl(b_i - a_i^\top x\bigr)
	=
	- \frac{a_i}{b_i - a_i^\top x}
	=
	- \frac{a_i}{s_i(x)}.
	\]
	Повторное дифференцирование даёт
	\[
	\nabla^2 \ln\bigl(b_i - a_i^\top x\bigr)
	=
	- \frac{a_i a_i^\top}{\bigl(b_i - a_i^\top x\bigr)^2}
	=
	- \frac{a_i a_i^\top}{s_i(x)^2}.
	\]
	Остаётся просуммировать по \(i=1,\dots,m\).
\end{proof}

\begin{Remark}{Квадратическая аппроксимация}{taylor}
	Во второй окрестности точки \(x_k\) функция \(f_k\) аппроксимируется формулой Тейлора второго порядка:
	\[
	\widetilde f_k(x)
	=
	f_k(x_k)
	+
	\nabla f_k(x_k)^\top (x - x_k)
	+
	\frac12 (x - x_k)^\top \nabla^2 f_k(x_k)\, (x - x_k).
	\]
	Точка максимума квадратической модели определяется из условия
	\[
	\nabla \widetilde f_k(x) = 0,
	\]
	то есть из линейной системы
	\[
	0 = \nabla f_k(x_k) + \nabla^2 f_k(x_k)\, (x - x_k).
	\]
	Отсюда получается шаг Ньютона
	\[
	x_{k+1}
	=
	x_k - \bigl(\nabla^2 f_k(x_k)\bigr)^{-1}\nabla f_k(x_k).
	\]
\end{Remark}

\subsubsection*{Эллипсоид Дикина}

\begin{Definition}{Эллипсоид Дикина}{dikin-ellipsoid}
	Для точки \(z\), удовлетворяющей \(A z < b\), и радиуса \(R > 0\) определим
	\[
	\cE(z,R)
	:=
	\left\{
	x \in \R^n :
	\sum_{i=1}^m
	\frac{\bigl(s_i(x)-s_i(z)\bigr)^2}{s_i(z)^2}
	\le R^2
	\right\}.
	\]
	Так как
	\[
	s_i(x)-s_i(z) = -a_i^\top (x-z),
	\]
	это равносильно записи
	\[
	\cE(z,R)
	=
	\left\{
	x \in \R^n :
	(x-z)^\top
	\left(
	\sum_{i=1}^m \frac{a_i a_i^\top}{s_i(z)^2}
	\right)
	(x-z)
	\le R^2
	\right\}.
	\]
\end{Definition}

\begin{Remark}{Назначение эллипсоида Дикина}{dikin-purpose}
	Эллипсоид Дикина задаёт естественную локальную геометрию барьерной функции и служит удобной метрикой
	для измерения близости точки \(x_k\) к центральной точке \(x_k^\ast\).
\end{Remark}

\subsubsection*{Схема алгоритма}

\begin{Statement}{Схема метода внутренней точки}{ipm-scheme}
	Пусть
	\[
	w_0 = 2^{-\Theta(L)},
	\qquad
	w_{k+1} = \left(1 + \frac{1}{100\sqrt m}\right) w_k.
	\]
	
	\begin{enumerate}
		\item Инициализация: выбрать внутреннюю точку \(x_0\).
		\item Для \(k=0,1,2,\dots\):
		\begin{enumerate}
			\item[] положить
			\[
			f_k(x)=w_k c^\top x+g(x);
			\]
			\item[] выполнить шаг Ньютона
			\[
			x_{k+1}
			=
			x_k-\bigl(\nabla^2 f_k(x_k)\bigr)^{-1}\nabla f_k(x_k);
			\]
			\item[] обновить параметр
			\[
			w_{k+1}
			=
			\left(1+\frac{1}{100\sqrt m}\right)w_k.
			\]
		\end{enumerate}
	\end{enumerate}
\end{Statement}

\begin{Statement}{Типовые инварианты и оценки}{ipm-estimates}
	В стандартном анализе метода внутренней точки доказываются следующие факты.
	
	\begin{enumerate}
		\item Если текущая точка удовлетворяет
		\[
		x_k \in \cE\!\left(x_k^\ast,\frac{1}{12}\right),
		\]
		то после умеренного увеличения параметра \(w_k\) новая центральная точка \(x_{k+1}^\ast\)
		остаётся близкой к \(x_k^\ast\).
		
		\item Один шаг Ньютона, выполненный из точки, достаточно близкой к \(x_k^\ast\),
		переносит итерацию в точку, которая достаточно близка к \(x_{k+1}^\ast\).
		
		\item В частности, можно поддерживать инвариант вида
		\[
		x_k \in \cE\!\left(x_k^\ast,\frac{1}{12}\right)
		\quad \text{для всех } k.
		\]
		
		\item Для оптимального значения \(x^\ast\) исходной задачи верна оценка разрыва
		\[
		\bigl|c^\top x_k - c^\top x^\ast\bigr| = O\!\left(\frac{m}{w_k}\right).
		\]
		
		\item Поэтому достаточно увеличить \(w_k\) до величины порядка
		\[
		w_k = \Theta\!\bigl(m\, 2^{\Theta(L)}\bigr),
		\]
		чтобы обеспечить точность
		\[
		\bigl|c^\top x_k - c^\top x^\ast\bigr| < 2^{-\Theta(L)}.
		\]
	\end{enumerate}
\end{Statement}

\begin{Proposition}{Число итераций}{iteration-bound}
	Пусть
	\[
	w_k = \left(1 + \frac{1}{100\sqrt m}\right)^k w_0,
	\qquad
	w_0 = 2^{-\Theta(L)}.
	\]
	Тогда для достижения точности \(2^{-\Theta(L)}\) достаточно
	\[
	k = O(\sqrt m\, L)
	\]
	итераций.
\end{Proposition}

\begin{proof}
	Необходимо довести \(w_k\) от величины \(2^{-\Theta(L)}\) до величины \(\Theta(m\,2^{\Theta(L)})\).
	Это эквивалентно условию
	\[
	\frac{w_k}{w_0} = m\, 2^{\Theta(L)}.
	\]
	Беря логарифм, получаем
	\[
	k \ln\!\left(1+\frac{1}{100\sqrt m}\right) = \Theta(L+\ln m).
	\]
	Так как \(\ln(1+\varepsilon)=\Theta(\varepsilon)\) при малых \(\varepsilon\), имеем
	\[
	\ln\!\left(1+\frac{1}{100\sqrt m}\right)=\Theta\!\left(\frac{1}{\sqrt m}\right).
	\]
	Следовательно,
	\[
	k = O\!\bigl(\sqrt m \, (L+\ln m)\bigr)=O(\sqrt m\,L).
	\]
\end{proof}

\begin{Example}{Единичный \(\ell_1\)-шар}{l1-ball}
	Множество
	\[
	\left\{
	x \in \R^n :
	|x_1|+\dots+|x_n| \le 1
	\right\}
	\]
	является выпуклым многогранником.
	
	\begin{enumerate}
		\item При \(n=2\) это ромб.
		\item При \(n=3\) это октаэдр.
	\end{enumerate}
	
	Именно такие фигуры были схематически изображены на доске как пример многогранников.
\end{Example}
